\documentclass[12pt]{article}

\usepackage[utf8]{inputenc}
\usepackage{amsmath, amssymb, amsthm}
\usepackage{geometry}
\geometry{a4paper, margin=2.5cm}

\title{Teorema Riesz--Schauder}
\author{}
\date{}

\theoremstyle{definition}
\newtheorem{definition}{Definiție}

\theoremstyle{plain}
\newtheorem{theorem}{Teoremă}
\newtheorem{proposition}{Propoziție}

\begin{document}

\maketitle

\begin{definition}
Fie $X$ un spațiu Banach. Un operator liniar continuu $T : X \to X$ se numește \emph{compact} dacă imaginea oricărei mulțimi mărginită este relativ compactă, adică


\[
B \subset X \text{ mărginită } \implies T(B) \text{ are aderență compactă}.
\]


Spectrul lui $T$ este mulțimea


\[
\sigma(T) = \{ \lambda \in \mathbb{C} : T - \lambda I \text{ nu este invertibil} \}.
\]


\end{definition}

\begin{theorem}[Riesz--Schauder]
Fie $X$ un spațiu Banach și $T : X \to X$ un operator compact. Atunci:
\begin{enumerate}
    \item[(i)] $0 \in \sigma(T)$.
    \item[(ii)] Orice $\lambda \in \sigma(T)$, $\lambda \ne 0$, este valoare proprie.
    \item[(iii)] Pentru $\lambda \ne 0$, spațiul propriu
    

\[
    E_\lambda = \ker(T - \lambda I)
    \]


    este finit dimensional.
    \item[(iv)] Singurele puncte de acumulare ale spectrului sunt $0$.
\end{enumerate}
\end{theorem}

\begin{proposition}
Pentru un operator compact $T : X \to X$, următoarele afirmații sunt echivalente pentru $\lambda \ne 0$:
\begin{enumerate}
    \item[(i)] $\lambda \in \sigma(T)$.
    \item[(ii)] $T - \lambda I$ nu este surjectiv.
    \item[(iii)] $T - \lambda I$ nu este injectiv.
    \item[(iv)] $\lambda$ este valoare proprie a lui $T$.
\end{enumerate}
\end{proposition}

\begin{proposition}
Fie $X$ un spațiu Banach și $T : X \to X$ un operator compact. Atunci pentru operatorul $I - T$ au loc următoarele proprietăți fundamentale:
\begin{enumerate}
    \item[(i)] Nucleul este finit dimensional,
    

\[
    N(I-T) = \ker(I-T) \quad \text{este finit dimensional};
    \]


    \item[(ii)] Imaginea este închisă,
    

\[
    R(I-T) = \operatorname{Im}(I-T) \quad \text{este închisă în } X;
    \]


    \item[(iii)] Următoarele afirmații sunt echivalente:
    

\[
    N(I-T) = \{0\}
    \quad \Longleftrightarrow \quad
    R(I-T) = X.
    \]


\end{enumerate}
\end{proposition}


\end{document}
